% Generated draft: requires mathematical and source review.
\documentclass[11pt]{article}
\usepackage[margin=25mm]{geometry}
\usepackage{fontspec}
\setmainfont{Latin Modern Roman}
\usepackage{amsmath,amssymb,mathtools,unicode-math}
\setmathfont{Latin Modern Math}
\usepackage{enumitem,needspace,xurl,hyperref}
\setlength{\parindent}{0pt}
\setlength{\parskip}{5pt}
\setlength{\emergencystretch}{3em}
\newcommand{\N}{\mathbb{N}}
\newcommand{\Z}{\mathbb{Z}}
\newcommand{\Q}{\mathbb{Q}}
\newcommand{\R}{\mathbb{R}}
\newcommand{\C}{\mathbb{C}}
\newcommand{\F}{\mathbb{F}}
\newcommand{\A}{\mathbb{A}}
\newcommand{\PP}{\mathbb{P}}
\newcommand{\E}{\mathbb{E}}
\newcommand{\eps}{\varepsilon}
\newcommand{\dd}{\,\mathrm{d}}
\begin{document}
\section*{UnsolvedMath ID: 20002871}
\textbf{Source problem code:} catalogue entry
\subsection{Definitions and mathematical statement}
Let
\[
p(z,\overline z,w,\overline w)=\sum_{\alpha,\beta,\gamma,\delta\geq0}
p_{\alpha\beta\gamma\delta}z^\alpha\overline z^\beta w^\gamma\overline w^\delta
\]
be a real-valued plurisubharmonic polynomial on $\C^2$, so its $2\times2$ complex Hessian is positive semidefinite at every point. Set
\[
N(p)=\{(\alpha+\beta,\gamma+\delta):p_{\alpha\beta\gamma\delta}\ne0\}.
\]
Use the lower-supporting-edge interpretation of an extreme Newton edge: a support line with no support points strictly below it, carrying the relevant lower face. In the usual positive-weight case the line has equation $ua+vb=d$ with $u,v>0$. Let $\Gamma_1,\Gamma_2$ be distinct extreme edges with
\[
\Gamma_1\cap\Gamma_2\ne\varnothing.
\]
The supplied source verification expressly resolves the corrupted OCR as nonempty intersection. Define the union truncation, counting any shared terms once, by
\[
q=\sum_{(\alpha+\beta,\gamma+\delta)\in\Gamma_1\cup\Gamma_2}
p_{\alpha\beta\gamma\delta}z^\alpha\overline z^\beta w^\gamma\overline w^\delta.
\]
B. Stensones asks whether $q=q_1+q_2$ with each $q_j$ real, plurisubharmonic and weighted homogeneous. Weighted homogeneity means $q_j(t^{u_j}z,t^{v_j}w)=t^{d_j}q_j(z,w)$ for $t>0$ and suitable weights $u_j,v_j>0$; the two summands may have different weights.

The original remark says $q$ is plurisubharmonic, although not necessarily weighted homogeneous. The aid reports that Simon--Stensones published a squared-norm plurisubharmonic polynomial for which the truncation on two neighboring extreme edges is not plurisubharmonic near the origin. If that assertion is verified, it simultaneously refutes the remark and prevents the requested decomposition, since a sum of plurisubharmonic polynomials is plurisubharmonic. No formula or full citation for this counterexample is supplied here. The aid contrasts genuinely disjoint edges, for which their separate weighted-homogeneous edge components give the decomposition.

A separate reduction in the aid concerns neighboring edges meeting at a single vertex $v=(a_0,b_0)$, each containing at least three distinct support points. Let $V$ be the entire component of bidegree $v$, and let $A$ and $B$ be the components exclusive to $\Gamma_1$ and $\Gamma_2$, respectively. Thus $q=A+B+V$. Under the usual positive-weight edge convention, the aid asserts that every two-weight decomposition is forced to use the original two edge lines, up to interchange. Its complete freedom is therefore a real polynomial $C$ supported at the single bidegree $v$:
\[
q_1=A+C,\qquad q_2=B+V-C.
\]
Writing $\mathcal L_f$ for the complex Hessian of $f$, existence of a plurisubharmonic decomposition in this case is equivalent to finding such a $C$ satisfying
\[
\mathcal L_A+\mathcal L_C\succeq0,\qquad
\mathcal L_B+\mathcal L_V-\mathcal L_C\succeq0
\quad\text{at every }(z,w)\in\C^2.
\]
This is finite-dimensional in the coefficients of $C$, with pointwise positivity constraints. The aid also reports a scalar-allocation obstruction for two transverse rank-one binomial edges sharing a positive modulus-square corner. Its exact edge polynomials and transversality hypotheses are omitted. Restricting to $C=tV$ is a special case and does not replace the arbitrary vertex-polynomial feasibility problem.
\subsection{Short English statement}
Ask whether truncation to two intersecting extreme Newton edges splits into two plurisubharmonic weighted-homogeneous pieces. The supplied aid reports a counterexample to even the truncation's plurisubharmonicity and gives a vertex-allocation criterion under an additional three-support-points-per-edge hypothesis.
\subsection{Sources}
\begin{itemize}
\item Local UnsolvedMath registry, permanent ID $20002871$, source code \texttt{catalogue entry}; canonical \texttt{aim-several-complex-variables-notes.json}, \texttt{notes[28]}.
\item AIM workshop \emph{Emerging applications of complexity for CR mappings}, Problem 2, attributed to B. Stensones, with its plurisubharmonicity remark: \url{https://aimath.org/WWN/crmappings/crmappings.pdf}. The supplied verification reports that the original PDF resolves OCR ``6 ='' as $\ne$ and specifies intersecting, not disjoint, edges. This draft relies on that supplied report, without independently accessing the PDF.
\item The aid attributes the negative example to Simon and Stensones, without title, publication date, explicit polynomial or authoritative full citation. The original literature field is empty. The trailing phrase ``CR manifolds and mappings'' is a section heading attached to the extracted remark.
\item Research attempt 1: \texttt{counterexample}, status at run \texttt{refuted}, current registry status \texttt{solved}, classification \texttt{SOLVED-IN-LITERATURE}, rigor \texttt{proved}, novelty confidence low. The corner-allocation reduction and literature assessment are machine-generated aid, not peer reviewed, and require independent verification before citation.
\end{itemize}
\end{document}
